% \documentclass[12pt]{article}

% \usepackage[margin=1in]{geometry}
% \usepackage{amsmath}
% \usepackage{amssymb}
% \usepackage{setspace}
% \usepackage{parskip}
% \usepackage{hyperref}
% \usepackage{booktabs}
% \usepackage{bibentry}

% \doublespacing

% \title{Methods: Spatial Distribution Analysis}
% \author{}
% \date{}

% \begin{document}

% \maketitle

% \section{Methods}

% \subsection{\quad Spatial Distribution Analysis via Morisita's Index of Dispersion}

% To quantify the spatial distribution of fluorescent particles in peptide assembly
% images, Morisita's Index of Dispersion ($I_M$) was computed using a quadrat-based
% random sampling approach. The complete analysis pipeline is described below.

% % -------------------------------------------------------------------
% \subsubsection{\quad Image Loading and Pre-processing}

% Fluorescence microscopy images acquired under varying buffer conditions
% (NaHEPES and Naphosphate, at concentrations of 10, 50, 100, and 250~mM, across
% pH values of 6.0, 6.5, 7.0, 7.5, and 8.0) were imported in grayscale format
% using the OpenCV library (\texttt{cv2.imread}, flag: \texttt{IMREAD\_GRAYSCALE}).
% Each image was represented as a two-dimensional array of pixel intensity values
% in the range $[0,\,255]$.

% % -------------------------------------------------------------------
% \subsubsection{\quad Image Binarization}

% Prior to spatial analysis, each grayscale image was converted to a binary mask to
% isolate fluorescent signal from background. Otsu's automatic thresholding
% method~\cite{otsu1979} was applied (\texttt{cv2.threshold} with
% \texttt{THRESH\_BINARY\,+\,THRESH\_OTSU}), which determines an optimal global
% intensity threshold by minimising the intra-class variance between foreground and
% background pixel populations. Pixels with intensity values exceeding the Otsu
% threshold were assigned a value of 1 (bright/signal), while all remaining pixels
% were assigned a value of 0 (background). This approach was adopted to ensure that
% binarization was objective and reproducible across all experimental conditions,
% without requiring manual threshold selection.

% % -------------------------------------------------------------------
% \subsubsection{\quad Random Quadrat Sampling}

% To assess spatial heterogeneity, $q = 200$ non-deterministic square quadrats of
% side length $s = 64$~pixels were randomly positioned within the binarized image.
% The top-left corner coordinates $(r_i,\,c_i)$ of each quadrat were drawn
% independently and uniformly from the set of valid positions --- defined as all
% positions for which the quadrat lay entirely within the image boundaries --- using
% a seeded pseudo-random number generator (NumPy \texttt{default\_rng}, seed $= 42$)
% to ensure reproducibility. For each quadrat $Q_i$, the number of bright pixels
% $n_i$ was recorded by summing all foreground pixels within the corresponding
% $s \times s$ pixel region of the binary mask:

% \begin{equation}
%     n_i \;=\; \sum_{(x,\,y)\,\in\, Q_i} B(x,\,y),
%     \qquad i = 1, 2, \ldots, q
%     \label{eq:quadrat_count}
% \end{equation}

% \noindent where $B(x,\,y) \in \{0,\,1\}$ denotes the binary mask value at pixel
% position $(x,\,y)$.

% % -------------------------------------------------------------------
% \subsubsection{\quad Computation of Morisita's Index of Dispersion}

% Morisita's Index of Dispersion ($I_M$) was computed from the quadrat counts
% following the formulation originally proposed by Morisita~\cite{morisita1959}:

% \begin{equation}
%     I_M \;=\; \frac{q\,\displaystyle\sum_{i=1}^{q} n_i\,(n_i - 1)}{N\,(N - 1)}
%     \label{eq:morisita}
% \end{equation}

% \noindent where $N = \sum_{i=1}^{q} n_i$ is the total number of bright pixels
% observed across all $q$ sampled quadrats. The index $I_M$ provides a
% dimensionless measure of spatial aggregation: a value of $I_M = 1$ is consistent
% with a random (Poisson) distribution of particles, $I_M > 1$ indicates spatial
% clustering, and $I_M < 1$ indicates a uniform or regular distribution. Images
% yielding $N \leq 1$ across all quadrats were excluded from analysis and flagged as
% containing insufficient fluorescent signal.

% % -------------------------------------------------------------------
% \subsubsection{\quad Statistical Significance Testing}

% The statistical significance of the observed $I_M$ value was assessed using the
% chi-squared test derived by Morisita~\cite{morisita1959}. Under the null hypothesis
% $H_0$ that the bright pixels are spatially randomly distributed ($I_M = 1$), the
% test statistic is given by:

% \begin{equation}
%     \chi^2 \;=\; I_M\,(N - 1) \;+\; q \;-\; N
%     \label{eq:chi2}
% \end{equation}

% \noindent with degrees of freedom $\mathrm{df} = q - 1$. The corresponding
% $p$-value was obtained from the chi-squared cumulative distribution function
% \cite{scipy2020}. A significance threshold of $\alpha = 0.05$ was applied:
% if $p < 0.05$, the null hypothesis was rejected and the distribution was
% classified as either \textit{clustered} ($I_M > 1$) or \textit{uniform}
% ($I_M < 1$); if $p \geq 0.05$, the distribution was classified as \textit{random}.

% % -------------------------------------------------------------------
% \subsubsection{\quad Visualisation}

% For each image, three diagnostic panels were generated:
% \begin{enumerate}
%     \item The original grayscale image with all sampled quadrats overlaid,
%           colour-coded by whether their bright-pixel count exceeded (red) or fell
%           below (cyan) the mean count $\bar{n} = N/q$ across all quadrats.
%     \item A frequency histogram of the per-quadrat bright-pixel counts $n_i$,
%           providing a visual assessment of count dispersion relative to a
%           Poisson-distributed baseline.
%     \item A bar chart of the computed $I_M$ value plotted against the $I_M = 1$
%           reference line, annotated with the chi-squared statistic and $p$-value.
% \end{enumerate}

% \noindent Results for all buffer conditions were compiled into a summary table and
% exported to a CSV file for downstream statistical comparison.

% % -------------------------------------------------------------------
% \bibliographystyle{ieeetr}
% \begin{thebibliography}{9}

% \bibitem{morisita1959}
% Morisita, M. (1959).
% \newblock Measuring the dispersion of individuals and analysis of the
% distributional patterns.
% \newblock \textit{Memoirs of the Faculty of Science, Kyushu University,
% Series~E (Biology)}, 2, 215--235.

% \bibitem{otsu1979}
% Otsu, N. (1979).
% \newblock A threshold selection method from gray-level histograms.
% \newblock \textit{IEEE Transactions on Systems, Man, and Cybernetics},
% 9(1), 62--66.

% \bibitem{scipy2020}
% Virtanen, P., Gommers, R., Oliphant, T.~E., et~al. (2020).
% \newblock {SciPy} 1.0: Fundamental algorithms for scientific computing in
% {Python}.
% \newblock \textit{Nature Methods}, 17, 261--272.

% \end{thebibliography}

% \end{document}



\documentclass[12pt]{article}

\usepackage[margin=1in]{geometry}
\usepackage{amsmath}
\usepackage{amssymb}
\usepackage{setspace}
\usepackage{parskip}
\usepackage{hyperref}
\usepackage{booktabs}
\usepackage{bibentry}
\usepackage{xcolor}

\doublespacing

\title{Methods: Spatial Distribution Analysis}
\author{}
\date{}

\begin{document}

\maketitle

\section{Methods: Spatial Distribution Analysis via Morisita's Index of Dispersion}

To quantify the spatial distribution of fluorescent particles in \textcolor{red}{membranes}, Morisita's Index of Dispersion ($I_M$) was computed using a quadrat-based
random sampling approach. The complete analysis pipeline is described below.

\begin{itemize}
    \item \textbf{Image Pre-processing :} Fluorescence microscopy images acquired under varying buffer conditions (NaHEPES and Naphosphate, at concentrations of 10, 50, 100, and 250~mM, across
    pH values of 6.0, 6.5, 7.0, 7.5, and 8.0) were imported in grayscale format
    using the OpenCV library.
    Each image was represented as a two-dimensional array of pixel intensity values
    in the range $[0,\,255]$.
    \item \textbf{Image Binarization :} Prior to spatial analysis, each grayscale image was converted to a binary mask to
    isolate fluorescent signal from background. Otsu's automatic thresholding
    method~\cite{otsu1979} was applied, which determines an optimal global
    intensity threshold by minimising the intra-class variance between foreground and
    background pixel populations. Pixels with intensity values exceeding the Otsu
    threshold were assigned a value of 1 (bright/signal), while all remaining pixels
    were assigned a value of 0 (background). This approach was adopted to ensure that
    binarization was objective and reproducible across all experimental conditions,
    without requiring manual threshold selection.

    \item \textbf{Random Quadrant Sampling :} To assess spatial heterogeneity, $q = 1000$ non-deterministic square quadrats of
    side length $s = 5$~pixels were randomly positioned within the binarized image. For each quadrat $Q_i$, the number of bright pixels
    $n_i$ was recorded by summing all foreground pixels within the corresponding
    $s \times s$ pixel region of the binarized image.

    \item \textbf{Computation of Morisita's Index of Dispersion :} Morisita's Index of Dispersion ($I_M$) was computed from the quadrat counts
    following the formulation originally proposed by Morisita~\cite{morisita1959}:
    
    \begin{equation}
        I_M \;=\; \frac{q\,\displaystyle\sum_{i=1}^{q} n_i\,(n_i - 1)}{N\,(N - 1)}
        \label{eq:morisita}
    \end{equation}
    
    \noindent where $N = \sum_{i=1}^{q} n_i$ is the total number of bright pixels
    observed across all $q$ sampled quadrats. The index $I_M$ provides a
    dimensionless measure of spatial aggregation, the 
    signification of which can be obtained using a hypothesis
    The testing framework is explained next. 

    \item \textbf{Statistical Significance Testing :} The statistical significance of the observed $I_M$ value was assessed using the
    chi-squared test derived by Morisita~\cite{morisita1959}. Under the null hypothesis
    $H_0$ that the bright pixels are spatially randomly distributed ($I_M = 1$), the
    test statistic is given by:
    
    \begin{equation}
        \chi^2 \;=\; I_M\,(N - 1) \;+\; q \;-\; N
        \label{eq:chi2}
    \end{equation}
    
    \noindent with degrees of freedom $\mathrm{df} = q - 1$. The corresponding
    $p$-value was obtained from the chi-squared cumulative distribution function
    \cite{scipy2020}. A significance threshold of $\alpha = 0.05$ was applied:
    if $p < 0.05$, the null hypothesis was rejected and the distribution was
    classified as either \textit{clustered} ($I_M > 1$) or \textit{uniform}
    ($I_M < 1$); if $p \geq 0.05$, the distribution was classified as \textit{random}.
    
\end{itemize}


% -------------------------------------------------------------------
\bibliographystyle{ieeetr}
\begin{thebibliography}{9}

\bibitem{morisita1959}
Morisita, M. (1959).
\newblock Measuring the dispersion of individuals and analysis of the
distributional patterns.
\newblock \textit{Memoirs of the Faculty of Science, Kyushu University,
Series~E (Biology)}, 2, 215--235.

\bibitem{otsu1979}
Otsu, N. (1979).
\newblock A threshold selection method from gray-level histograms.
\newblock \textit{IEEE Transactions on Systems, Man, and Cybernetics},
9(1), 62--66.

\bibitem{scipy2020}
Virtanen, P., Gommers, R., Oliphant, T.~E., et~al. (2020).
\newblock {SciPy} 1.0: Fundamental algorithms for scientific computing in
{Python}.
\newblock \textit{Nature Methods}, 17, 261--272.

\end{thebibliography}

\end{document}

